
@@M0@@ \uline{soit exacte).}

Posons @@M1@@, @@M2@@ et définissons par récurrence @@M3@@. On vérifie aussitôt que @@M4@@. Reste à prouver que pour @@M5@@ assez grand, on a @@M6@@. En effet considérons la suite décroissante de @@M7@@-modules @@M8@@ et soit @@M9@@. La suite décroissante de fermés @@M10@@ de @@M11@@ est stationnaire, soit @@M12@@ la valeur stationnaire. Montrons que @@M13@@. Sinon soit @@M14@@ un point maximal de @@M15@@. Quitte à remplacer @@M16@@ par @@M17@@, on peut supposer @@M18@@ local de point fermé @@M19@@. Vu le choix de @@M20@@, @@M21@@ est artinien pour @@M22@@ assez grand, donc la suite @@M23@@ est stationnaire. Il existe donc @@M24@@ tel que @@M25@@. Par suite @@M26@@ est un monomorphisme, donc une immersion fermée (Nakayama); comme @@M27@@, on a @@M28@@, et @@M29@@, absurde.

\section*{3. Démonstration de I : le dévissage de Oort}

\textbf{Proposition 3.1.} (Oort, Bull. soc. math. France, t. 90-1962). \uline{Soient} @@M30@@ \uline{un schéma localement noethérien,} @@M31@@ \uline{un morphisme propre,} @@M32@@ \uline{un sous-schéma fermé. On suppose que~:}
\begin{enumerate}[label=\roman*)]
\item @@M33@@ \uline{est défini par un idéal} @@M34@@ \uline{contenu dans le nilradical} @@M35@@ \uline{de} @@M36@@ \uline{et} @@M37@@.
\item @@M38@@ \uline{est plat sur} @@M39@@ \uline{et cohomologiquement plat sur} @@M40@@ \uline{en dimension} @@M41@@ (EGA III 7.8.1).
\item @@M42@@ \uline{est plat sur} @@M43@@ \uline{et cohomologiquement plat sur} @@M44@@ \uline{en dimension} @@M45@@ \uline{et} @@M46@@ \uline{de sorte que (d'après} EGA III 7.4.6\uline{) il existe un} @@M47@@\uline{-module cohérent} @@M48@@ \uline{et un isomorphisme fonctoriel en le} @@M49@@\uline{-module cohérent} @@M50@@,
